\chapter{Assembly as a Fold}
\label{ch:assembly}
\label{ch:fem}% alias: "the finite-element discretization/assembly" is cited as ch:fem in later chapters

\chapterorient{Chapter~\ref{ch:galerkin} told us \emph{what} the global matrix
is---the bilinear form sampled on basis pairs---and Chapter~\ref{ch:refelement}
told us how to compute one element's worth of it on a reference cell. This
chapter is about the \emph{seam} between the two: how the small, dense
element matrices, each living in its own local coordinate system, are stitched
into the one big sparse global matrix the solver actually sees. The mechanism
is a pair of indexing operations---\emph{gather} and \emph{scatter}---and the
moment we name them correctly, a structural fact about the whole simulator
comes into focus. Assembly is not a tangle of nested loops; it is a
\emph{fold}: a list of contributions combined, one at a time, by a single
associative operation into a single result. That sentence is the seed of an
idea this book spends Part~III harvesting---that the simulator's deep structure
is folds over lists. By the end you will have assembled the 1-D Poisson matrix
by hand a second time, but now the \emph{right} way, watching shared nodes
collect contributions from both their neighbours, and you will have proved that
the order in which terms and elements are accumulated never matters.}

\section{The gap between local and global}

We left Chapter~\ref{ch:galerkin} with a clean formula and a quiet omission. The
formula was $\mat{K}_{ij}=a(\phi_j,\phi_i)$: every entry of the global stiffness
matrix is the bilinear form evaluated on a pair of global basis functions. The
omission was \emph{how you actually compute it}. Written as one integral over
the whole domain $\dom$, the entry
\begin{equation}
\mat{K}_{ij}=\int_\dom Q\,\diffop\phi_j\cdot\diffop\phi_i\dV
\label{eq:global-entry}
\end{equation}
is honest but useless as a recipe: the global basis functions $\phi_i$ are
awkward piecewise objects, defined cell by cell, and nobody integrates them as
single functions over the entire mesh. Real finite-element codes never form
\cref{eq:global-entry} directly. They do something that looks, at first, like a
detour: they compute small \emph{element matrices}, one per mesh cell, each
holding only the interactions among the handful of basis functions that live on
that one cell, and then they \emph{add} those small matrices into the big one.
This chapter is about why that detour is in fact the royal road, and about the
algebraic shape it gives to the whole assembly stage.

The key observation is that the domain integral \cref{eq:global-entry} splits
across cells. The mesh $\mathcal{T}=\{K_1,\dots,K_{n_e}\}$ partitions $\dom$ into
$n_e$ non-overlapping elements, so any integral over $\dom$ is the sum of
integrals over the elements:
\begin{equation}
\mat{K}_{ij}=\int_\dom Q\,\diffop\phi_j\cdot\diffop\phi_i\dV
=\sum_{K\in\mathcal{T}}\;\int_K Q\,\diffop\phi_j\cdot\diffop\phi_i\dV.
\label{eq:cellsplit}
\end{equation}
Already \cref{eq:cellsplit} has a summation in it---our first sum over elements,
and the structural heart of the chapter. But most of its terms are zero. On a
given cell $K$, the only basis functions that are nonzero are the few whose
support touches $K$; for the hat basis of \cref{ch:galerkin} that is exactly the
basis functions sitting on the \emph{vertices of $K$}. If either $\phi_i$ or
$\phi_j$ vanishes on all of $K$, the cell's contribution to $\mat{K}_{ij}$ is
zero. So each cell contributes to only a tiny block of the global matrix: the
rows and columns indexed by \emph{its own} degrees of freedom.

\begin{keyidea}
Assembly inverts the two sums. \Cref{eq:cellsplit} loops over the global entries
$(i,j)$ and, for each, sums over cells---most contributing nothing. A
finite-element code loops the other way: \emph{over cells} on the outside, and
for each cell computes the small dense block of interactions among only that
cell's degrees of freedom, then \emph{adds} that block into the global matrix at
the right rows and columns. The same number $\sum_K\int_K(\cdots)$ is computed,
but the loop that is cheap (over elements, with dense local work) is on the
outside, and the loop that is expensive (over the $N^2$ global entries, almost
all zero) never happens.
\end{keyidea}

To make ``adds that block in at the right rows and columns'' precise we need a
language for moving between the global degrees of freedom and an element's local
ones. That language is gather and scatter.

\section{Gather and scatter: the operator \texorpdfstring{$G$}{G} and its transpose}
\label{sec:gatherscatter}

Number the global degrees of freedom $1,\dots,N$; for the hat basis these are
the interior mesh nodes. Each element $K$ owns a short list of \emph{local}
degrees of freedom---for a 1-D linear element, its two endpoints; for a triangle,
its three vertices---and there is a fixed table, the \emph{local-to-global map},
that says which global number each local slot corresponds to. This table is the
only bookkeeping assembly needs, and it is built once from the mesh.

\begin{definition}[Element gather]
For an element $K$ with $L$ local degrees of freedom, the \emph{gather} (or
\emph{restriction}) operator $\mat{R}_K$ is the $L\times N$ matrix that copies the
global degrees of freedom belonging to $K$ into a length-$L$ local vector, in the
element's local order:
\[
(\mat{R}_K\vx)_{\,\ell}=\vx_{\,g(K,\ell)},\qquad \ell=1,\dots,L,
\]
where $g(K,\ell)$ is the global index of $K$'s $\ell$-th local degree of freedom.
Each row of $\mat{R}_K$ is a single $1$ in column $g(K,\ell)$ and zeros
elsewhere: $\mat{R}_K$ is a \emph{Boolean selection matrix}, carrying no
arithmetic.
\end{definition}

The gather \emph{reads}. Its transpose \emph{writes}, and writes in a particular,
crucial way.

\begin{definition}[Element scatter-add]
The \emph{scatter-add} (or \emph{assembly}) operator is the transpose
$\mat{R}_K^{\mathsf T}$, an $N\times L$ matrix. Applied to a length-$L$ local
vector $\vu$, it places each local entry into its global slot and
\emph{accumulates}:
\[
(\mat{R}_K^{\mathsf T}\vu)_{\,m}=\sum_{\ell\,:\,g(K,\ell)=m}\vu_{\,\ell}.
\]
For a single element each global slot receives at most one local entry, but when
the scatter-adds of \emph{many} elements are summed---which is exactly what
assembly does---a global degree of freedom shared by several elements receives a
contribution from each.
\end{definition}

\Cref{fig:gatherscatter} draws the pair. Gather pulls a copy of element $K$'s
nodal values out of the long global vector; element-local computation happens on
that short vector; scatter-add pushes the results back, summing into any slot two
elements share.

\begin{figure}[t]
\centering
\begin{tikzpicture}[scale=1.0, every node/.style={font=\footnotesize}]
  % ---- global vector (vertical stack of cells) ----
  \def\cw{0.62}
  \foreach \k/\lab in {0/$x_1$,1/$x_2$,2/$x_3$,3/$x_4$,4/$x_5$}{
    \draw[simInk] (0,-\k*\cw) rectangle (\cw,-\k*\cw-\cw);
    \node at (\cw/2,-\k*\cw-\cw/2) {\lab};
  }
  \node[above, font=\footnotesize\bfseries] at (\cw/2,0.05) {global $\vx$};
  % highlight the two dofs of element K2: x2, x3 (indices 1,2)
  \draw[simBlue, very thick] (0,-1*\cw) rectangle (\cw,-3*\cw);
  % ---- element-local vector for K2 ----
  \def\ex{4.4}
  \foreach \k/\lab in {0/$u_1^{K_2}$,1/$u_2^{K_2}$}{
    \draw[simBlue, fill=simBgBlue] (\ex,-\k*\cw-0.62) rectangle (\ex+\cw,-\k*\cw-0.62-\cw);
    \node at (\ex+\cw/2,-\k*\cw-0.62-\cw/2) {\lab};
  }
  \node[above, font=\footnotesize\bfseries, simBlue!70!simInk] at (\ex+\cw/2,-0.5) {element $K_2$};
  % gather arrow
  \draw[flow, simBlue] (\cw+0.1,-1.4*\cw) to[bend left=12]
     node[above, font=\scriptsize]{$\mat{R}_{K_2}$ (gather)} (\ex-0.1,-1.0);
  % element computation box
  \def\bx{4.2}
  \node[stage, minimum width=20mm] (comp) at (\ex+\cw/2,-3.4) {compute\\$\mat{K}^{\mathrm{elem}}_{K_2}\vu$};
  \draw[flow, simBlue] (\ex+\cw/2,-2.0) -- (comp.north);
  % scatter-add arrow back, into x2 and x3
  \draw[backflow] (comp.west) to[bend right=18]
     node[below, font=\scriptsize, align=center]{$\mat{R}_{K_2}^{\mathsf T}$\\(scatter-add)} (\cw+0.1,-2.6*\cw);
  % a SECOND element K1 also writing into x2 (the shared dof)
  \draw[backflow, simGold] (-1.5,-1.6*\cw) to[bend left=20]
     node[left, font=\scriptsize, align=center]{from\\$K_1$} (-0.05,-1.5*\cw);
  \node[simGold!70!simInk, font=\scriptsize, align=center] at (-1.85,-3.0*\cw)
    {$x_2$ shared:\\two contributions};
  % bracket the shared dof
  \draw[simGold, thick] (-0.06,-1*\cw) -- (-0.06,-2*\cw);
\end{tikzpicture}
\caption[Gather and scatter-add]{The gather/scatter picture, for element $K_2$ of
a 1-D mesh. \emph{Gather} ($\mat{R}_{K_2}$) copies the two global degrees of
freedom belonging to $K_2$---here $x_2$ and $x_3$---into a short element-local
vector. The element matrix acts on that local vector. \emph{Scatter-add}
($\mat{R}_{K_2}^{\mathsf T}$) writes the results back into the global slots,
\emph{accumulating} rather than overwriting. The global degree of freedom $x_2$
sits on the boundary between $K_1$ and $K_2$, so it receives one contribution
from each element (the gold arrow is $K_1$'s)---this superposition is what
\emph{assembly} means.}
\label{fig:gatherscatter}
\end{figure}

\subsection{Assembly multiplicity and the diagonal \texorpdfstring{$\mat{R}^{\mathsf T}\mat{R}$}{R'R}}

The fact that worries newcomers---a shared node being written by two
elements---is not a bug to be guarded against; it \emph{is} the discretization.
A basis function $\phi_i$ on a shared node is, by construction, supported on
\emph{both} adjacent cells, so the integral $\int_\dom(\cdots)$ that defines its
matrix row genuinely receives a piece from each cell (\cref{eq:cellsplit}). The
scatter-\emph{add} reproduces that sum exactly. Overwriting instead of adding
would silently drop half of each shared row.

There is a clean way to measure how much sharing a mesh has. Stack the gathers of
all elements and ask: what does ``gather everything, then immediately
scatter-add it all back'' do to the global vector? For one element,
$\mat{R}_K^{\mathsf T}\mat{R}_K$ is the diagonal matrix that is $1$ on every
global slot $K$ touches and $0$ elsewhere. Summing over all elements,
\begin{equation}
\sum_{K\in\mathcal{T}}\mat{R}_K^{\mathsf T}\mat{R}_K
=\operatorname{diag}(\,m_1,m_2,\dots,m_N\,),
\qquad
m_i=\#\{\,K:\text{ global dof }i\text{ touches }K\,\},
\label{eq:multiplicity}
\end{equation}
the diagonal matrix whose $i$-th entry $m_i$ is the \emph{multiplicity} of global
degree of freedom $i$---the number of elements sharing it. For an interior 1-D
node, $m_i=2$; for a boundary node, $m_i=1$; for an interior vertex of a 2-D
triangular mesh, $m_i$ is the number of triangles in the fan around it.

\begin{notationbox}
The single-element operators are written $\mat{R}_K$ throughout, matching the
\code{element\_restrict} operator of the synthesis surface: $\mat{R}_K$ is the
``gather'' $G$ and $\mat{R}_K^{\mathsf T}$ is the ``scatter-add'' $G^{\mathsf T}$.
The diagonal \cref{eq:multiplicity} is the \emph{dof-multiplicity diagonal}
$G^{\mathsf T}G$. Note carefully: $\mat{R}_K^{\mathsf T}\mat{R}_K$ is a diagonal of
ones and zeros (a single element selects each of its dofs once), but the
\emph{sum over elements} \cref{eq:multiplicity} counts multiplicities. The
companion product $G\,G^{\mathsf T}$ on the element-local side is \emph{not} the
identity---it averages a shared value and redistributes it---a fact worth keeping
straight when you later meet it in the matrix-free apply of \cref{ch:matrixfree}.
\end{notationbox}

\section{Element matrices to global matrix}
\label{sec:elem-to-global}

With gather and scatter named, the assembly formula writes itself. On element
$K$, compute the $L\times L$ \emph{element matrix}
\begin{equation}
\bigl(\mat{K}^{\mathrm{elem}}_K\bigr)_{\ell\ell'}
=\int_K Q\,\diffop\psi^K_{\ell'}\cdot\diffop\psi^K_{\ell}\dV,
\label{eq:elemmat}
\end{equation}
where $\psi^K_1,\dots,\psi^K_L$ are the \emph{local} basis functions on $K$ (the
restrictions to $K$ of the global hats touching it, expressed in $K$'s own
coordinates---this is the reference-element computation of
\cref{ch:refelement}). The element matrix is small and dense: $2\times2$ for a
1-D linear element, $3\times3$ for a linear triangle, a few hundred entries even
for high-order 3-D elements. Now scatter each element matrix into the global one
and sum:
\begin{equation}
\boxed{\;\mat{K}=\sum_{K\in\mathcal{T}}\mat{R}_K^{\mathsf T}\,
\mat{K}^{\mathrm{elem}}_K\,\mat{R}_K\;}
\label{eq:assembly}
\end{equation}
The sandwich $\mat{R}_K^{\mathsf T}(\cdots)\mat{R}_K$ does exactly what its two
factors say: $\mat{R}_K$ on the right gathers the global indices into local order
for the columns, $\mat{R}_K^{\mathsf T}$ on the left scatters the local rows back
to global order, and the surrounding sum accumulates the overlapping blocks. That
\cref{eq:assembly} equals the entrywise formula \cref{eq:cellsplit} is immediate:
the $(i,j)$ entry of $\mat{R}_K^{\mathsf T}\mat{K}^{\mathrm{elem}}_K\mat{R}_K$ is
$(\mat{K}^{\mathrm{elem}}_K)_{\ell\ell'}$ when $i=g(K,\ell)$ and $j=g(K,\ell')$
for local slots $\ell,\ell'$ of $K$, and zero otherwise---precisely the cell-$K$
term of \cref{eq:cellsplit}, written in local coordinates.

\begin{figure}[t]
\centering
\begin{tikzpicture}[scale=0.66, every node/.style={font=\scriptsize}]
  % global 5x5 grid
  \def\N{5}
  \foreach \r in {0,...,4}{\foreach \c in {0,...,4}{
      \draw[simGray!35] (\c,-\r) rectangle ++(0.96,-0.96);}}
  \draw[simInk, thick] (0,0) rectangle (\N,-\N);
  % element K1 block: dofs {0,1} -> overlapping diagonal block rows/cols 0,1
  \fill[simBlue!35] (0,0) rectangle ++(1.96,-1.96);
  % element K2 block: dofs {1,2}
  \fill[simRust!35] (1,-1) rectangle ++(1.96,-1.96);
  % element K3 block: dofs {2,3}
  \fill[simGreen!35] (2,-2) rectangle ++(1.96,-1.96);
  % element K4 block: dofs {3,4}
  \fill[simViolet!35] (3,-3) rectangle ++(1.96,-1.96);
  % the OVERLAP cells (shared dofs) get two colors stacked: mark with hatch
  \foreach \d in {1,2,3}{
    \draw[simInk, thick] (\d,-\d) rectangle ++(0.96,-0.96);
    \node[font=\tiny\bfseries] at (\d+0.48,-\d-0.48) {$+$};
  }
  % labels
  \node[simBlue!70!simInk] at (0.96,0.45) {$K_1$};
  \node[simRust!70!simInk] at (3.4,-1.0) {$K_2$};
  \node[simGreen!70!simInk] at (4.5,-2.4) {$K_3$};
  \node[simViolet!70!simInk] at (4.7,-3.7) {$K_4$};
  \node[anchor=west, align=left, text width=30mm] at (5.5,-1.2)
    {Each element's $2\times2$ block lands on a $2\times2$ diagonal patch.};
  \node[anchor=west, align=left, text width=30mm] at (5.5,-3.1)
    {Shared dofs (marked $+$) sit where two blocks \emph{overlap}: the entries \emph{add}.};
\end{tikzpicture}
\caption[Scatter-add building the global matrix]{Assembly of a 1-D mesh with four
elements $K_1,\dots,K_4$ on five nodes (essential boundaries dropped for
clarity). Each element's small $2\times2$ matrix scatters onto the $2\times2$
diagonal patch indexed by its two nodes. Consecutive elements share one node, so
consecutive patches \emph{overlap} on a single diagonal entry (marked $+$): there
the two element contributions are \emph{added}. The overlapping-blocks-on-the-
diagonal picture is the geometric content of $\mat{K}=\sum_K\mat{R}_K^{\mathsf T}
\mat{K}^{\mathrm{elem}}_K\mat{R}_K$, and it is why the assembled matrix comes out
banded.}
\label{fig:scatteradd}
\end{figure}

\begin{workedexample}{Assembling 1-D Poisson by gather/scatter}{assemble-1d}
\label{wex:assemble-1d}
We re-derive the tridiagonal stiffness of \cref{ch:galerkin}, but this time the
\emph{assembly} way---element by element, scatter-adding---so the multiplicity on
shared nodes is explicit. Take the mesh of $n=4$ cells on $[0,1]$, $h=\tfrac14$,
with nodes $x_0,\dots,x_4$. The interior nodes $x_1,x_2,x_3$ carry the global
degrees of freedom $1,2,3$; the boundary nodes carry none. The four elements and
the global indices of their two local slots are
\[
\begin{array}{c|c|c}
\text{element} & \text{local slot }1 & \text{local slot }2\\\hline
K_1=[x_0,x_1] & (\text{bdry}) & 1\\
K_2=[x_1,x_2] & 1 & 2\\
K_3=[x_2,x_3] & 2 & 3\\
K_4=[x_3,x_4] & 3 & (\text{bdry})
\end{array}
\]

\emph{The element matrix.} Every linear element here is identical. With local
hats of slope $\pm1/h$ on a cell of length $h$, the $2\times2$ element stiffness
$\bigl(\mat{K}^{\mathrm{elem}}\bigr)_{\ell\ell'}=\int_K(\psi_{\ell'})'(\psi_\ell)'$
is
\[
\mat{K}^{\mathrm{elem}}=\frac1h\begin{pmatrix}\phantom{-}1 & -1\\ -1 & \phantom{-}1\end{pmatrix}
=4\begin{pmatrix}\phantom{-}1 & -1\\ -1 & \phantom{-}1\end{pmatrix}\quad(h=\tfrac14).
\]
The diagonal $+1/h$ is one node's $(\text{slope})^2\cdot h=(1/h)^2 h$; the
off-diagonal $-1/h$ is the two opposing slopes $(-1/h)(+1/h)\cdot h$.

\emph{Scatter, element by element.} Start with $\mat{K}=\mathbf 0$ (the
empty-mesh operator). Now add each $\mat{R}_K^{\mathsf T}\mat{K}^{\mathrm{elem}}
\mat{R}_K$, dropping rows/columns that scatter to a boundary node.
\begin{itemize}[nosep]
\item $K_1$ touches only global dof $1$ (its other slot is boundary). It
contributes $+\tfrac1h$ to $\mat{K}_{11}$.
\item $K_2$ touches dofs $1,2$. It contributes $\tfrac1h\!\begin{psmallmatrix}1&-1\\-1&1\end{psmallmatrix}$
to the $\{1,2\}$ block: $+\tfrac1h$ to each of $\mat{K}_{11},\mat{K}_{22}$ and
$-\tfrac1h$ to $\mat{K}_{12},\mat{K}_{21}$.
\item $K_3$ touches dofs $2,3$: $+\tfrac1h$ to $\mat{K}_{22},\mat{K}_{33}$,
$-\tfrac1h$ to $\mat{K}_{23},\mat{K}_{32}$.
\item $K_4$ touches only dof $3$: $+\tfrac1h$ to $\mat{K}_{33}$.
\end{itemize}
Now read the diagonal as a running tally. Global dof $1$ collected $+\tfrac1h$
from $K_1$ \emph{and} $+\tfrac1h$ from $K_2$: total $\tfrac2h$. Dof $2$ collected
$+\tfrac1h$ from $K_2$ and $+\tfrac1h$ from $K_3$: total $\tfrac2h$. Dof $3$,
symmetrically, $\tfrac2h$. \emph{Each diagonal entry is $\tfrac2h$ because each
interior node is shared by exactly two elements}---multiplicity $m_i=2$, exactly
\cref{eq:multiplicity}. The off-diagonals each got a single $-\tfrac1h$ from the
one element straddling that pair. Assembling,
\[
\mat{K}=\frac1h
\begin{pmatrix} 2 & -1 & 0\\ -1 & 2 & -1\\ 0 & -1 & 2\end{pmatrix},
\]
the tridiagonal matrix \cref{eq:tridiag}, now \emph{earned} as a superposition of
four $2\times2$ blocks rather than read off a stencil. The ``$2$'' on the
diagonal is not a coincidence of the second-derivative formula---it is two
elements each depositing a $1$.
\end{workedexample}

\subsection{Why this is the right factoring}

Equation~\cref{eq:assembly} cleanly separates three concerns that a single global
integral would tangle. The \emph{geometry and physics} live entirely inside
$\mat{K}^{\mathrm{elem}}_K$---the material weight $Q$, the element's shape, the
choice of basis. The \emph{connectivity} lives entirely inside the Boolean
$\mat{R}_K$---which local slot is which global node. And the \emph{combination}
is a plain sum. Each can be changed without touching the others: refine the mesh
and only the $\mat{R}_K$ table grows; change the material and only the element
matrices change; add a physical effect and---as the next section shows---only the
list of summands lengthens. This separation is what lets one assembly routine
serve every analysis in the book.

\begin{figure}[t]
\centering
\begin{tikzpicture}[scale=1.0, every node/.style={font=\footnotesize}]
  % two triangles sharing edge 2-3
  \coordinate (n1) at (0,1.4);
  \coordinate (n2) at (1.4,2.0);
  \coordinate (n3) at (1.0,0.0);
  \coordinate (n4) at (2.6,0.9);
  \filldraw[simBgBlue, draw=simBlue, thick] (n1)--(n2)--(n3)--cycle;
  \filldraw[simBgRust, draw=simRust, thick, fill opacity=0.6] (n2)--(n3)--(n4)--cycle;
  \draw[simGold, very thick] (n2)--(n3);
  \foreach \p/\lab/\pos in {n1/1/left,n2/2/above,n3/3/below,n4/4/right}{
    \fill[simInk] (\p) circle (2pt);
    \node[\pos=2pt] at (\p) {$\lab$};}
  \node[simBlue!70!simInk] at (0.8,1.25) {$K_a$};
  \node[simRust!70!simInk] at (1.8,0.95) {$K_b$};
  \node[simGold!75!simInk, font=\scriptsize, right] at (1.5,1.0) {shared edge $2$--$3$};
\end{tikzpicture}
\caption[Two triangles sharing an edge]{The two-triangle mesh of
\cref{wex:assemble-2d}. Triangles $K_a=\{1,2,3\}$ and $K_b=\{2,3,4\}$ share the
edge joining nodes $2$ and $3$ (gold). The shared nodes $2,3$ have multiplicity
$2$; the unshared nodes $1,4$ have multiplicity $1$. Each triangle's $3\times3$
element matrix scatters onto its three nodes, and the two scatters overlap on the
$2\times2$ block indexed by $\{2,3\}$---the discrete image of the shared edge.}
\label{fig:twotri}
\end{figure}

\begin{workedexample}{Two triangles, a shared edge, multiplicity in 2-D}{assemble-2d}
\label{wex:assemble-2d}
The 1-D example made every interior node have multiplicity exactly $2$. To see
the gather/scatter mechanism cope with richer connectivity---and to watch a shared
\emph{edge}, not just a shared node---we assemble on the smallest interesting 2-D
mesh: two linear triangles glued along an edge (\cref{fig:twotri}). Label the four
nodes $1,2,3,4$; triangle $K_a$ has vertices $\{1,2,3\}$ and triangle $K_b$ has
vertices $\{2,3,4\}$. Nodes $2$ and $3$ are shared (they lie on the common edge);
nodes $1$ and $4$ belong to one triangle each.

\emph{The gather tables.} In local vertex order,
\[
\mat{R}_{K_a}:\ (\ell_1,\ell_2,\ell_3)\mapsto(1,2,3),
\qquad
\mat{R}_{K_b}:\ (\ell_1,\ell_2,\ell_3)\mapsto(2,3,4).
\]
Write each triangle's $3\times3$ element stiffness abstractly as
$\mat{K}^{\mathrm{elem}}_{K}=\bigl(s^K_{pq}\bigr)_{p,q=1}^3$ (its actual entries
come from the reference-triangle computation of \cref{ch:refelement}; we need only
\emph{where} they land). Scatter-adding both into the $4\times4$ global matrix:
\[
\mat{K}=\mat{R}_{K_a}^{\mathsf T}\mat{K}^{\mathrm{elem}}_{K_a}\mat{R}_{K_a}
       +\mat{R}_{K_b}^{\mathsf T}\mat{K}^{\mathrm{elem}}_{K_b}\mat{R}_{K_b}
=\begin{pmatrix}
s^a_{11} & s^a_{12} & s^a_{13} & 0\\[1pt]
s^a_{21} & s^a_{22}\!+\!s^b_{11} & s^a_{23}\!+\!s^b_{12} & s^b_{13}\\[1pt]
s^a_{31} & s^a_{32}\!+\!s^b_{21} & s^a_{33}\!+\!s^b_{22} & s^b_{23}\\[1pt]
0 & s^b_{31} & s^b_{32} & s^b_{33}
\end{pmatrix}.
\]
Read the four corners of the sharing off the matrix. The $(1,4)$ and $(4,1)$
entries are \emph{zero}: nodes $1$ and $4$ never share a triangle, so no element
deposits there---exactly the sparsity rule of \cref{sec:cost}. The
$\{2,3\}\times\{2,3\}$ block is where the two scatters \emph{overlap}, and every
entry of it is a \emph{sum} $s^a_{\cdot\cdot}+s^b_{\cdot\cdot}$ of a contribution
from each triangle: the discrete image of the shared edge. The diagonal entries
$\mat{K}_{22}=s^a_{22}+s^b_{11}$ and $\mat{K}_{33}=s^a_{33}+s^b_{22}$ are two-term
sums (multiplicity $2$), while $\mat{K}_{11}=s^a_{11}$ and $\mat{K}_{44}=s^b_{33}$
are single terms (multiplicity $1$)---precisely
$\sum_K\mat{R}_K^{\mathsf T}\mat{R}_K=\operatorname{diag}(1,2,2,1)$ of
\cref{eq:multiplicity}. Nothing about the mechanism changed from 1-D; only the
element matrices grew from $2\times2$ to $3\times3$ and the gather tables grew a
column. The fold is identical.
\end{workedexample}

\section{Assembly is a fold over weak-form terms}
\label{sec:fold}

So far we have assembled \emph{one} operator from \emph{one} bilinear form, summed
over elements. Real analyses ask for more. The driven Maxwell operator is the sum
of three forms, $\mat{A}(\omega)=\mat{K}+i\omega\mat{C}-\omega^2\mat{M}$; a lossy
or multi-material problem adds still more terms; the mass matrix and the stiffness
matrix are two \emph{terms} an eigenmode analysis needs side by side. The global
operator an analysis assembles is a sum over a \emph{list of weak-form terms},
and each term's own assembly is, as we have just seen, a sum over elements. Two
sums, nested---and both are folds.

\subsection{The term list and its per-term operator}

Recall from \cref{sec:family} that every term of the bilinear-form family has the
shape $a_t(u,v)=(Q_t\,\diffop_t u,\diffop_t v)$: a material weight $Q_t$ and a
differential operator $\diffop_t\in\{\Grad,\,\mathrm{id},\,\Curl,\,\Div\}$.
Sampling one term on basis pairs and assembling it over the mesh
(\cref{eq:assembly}) produces one global matrix; call it $\mat{A}_t$, the
\emph{per-term operator}. An analysis hands the assembler a list of terms
$[\,t_1,t_2,\dots,t_m\,]$ and wants the matrix of the total form $a=\sum_t a_t$.
Because sampling and assembly are linear in the form,
\begin{equation}
\mat{K}=\sum_{t}\mat{A}_t,
\qquad
\mat{A}_t=\sum_{K\in\mathcal{T}}\mat{R}_K^{\mathsf T}\,
\mat{A}^{\mathrm{elem}}_{t,K}\,\mat{R}_K .
\label{eq:termsum}
\end{equation}
The outer sum is over \emph{terms}; the inner sum, hidden inside each
$\mat{A}_t$, is over \emph{elements}. \Cref{fig:termfold} draws the outer one:
each weak-form term enters its own assembly and the per-term operators are summed
into the global operator.

\begin{figure}[t]
\centering
\begin{tikzpicture}[node distance=5mm and 12mm, every node/.style={font=\footnotesize}]
  \node[opbox] (t1) {$t_1=(\varepsilon,\Grad)$};
  \node[opbox, below=of t1] (t2) {$t_2=(\sigma,\mathrm{id})$};
  \node[opbox, below=of t2] (t3) {$t_3=(\nu,\Curl)$};
  \node[stage, right=16mm of t1] (a1) {assemble};
  \node[stage, right=16mm of t2] (a2) {assemble};
  \node[stage, right=16mm of t3] (a3) {assemble};
  \draw[flow] (t1) -- (a1);
  \draw[flow] (t2) -- (a2);
  \draw[flow] (t3) -- (a3);
  \node[right=4mm of a1] (m1) {$\mat{A}_1$};
  \node[right=4mm of a2] (m2) {$\mat{A}_2$};
  \node[right=4mm of a3] (m3) {$\mat{A}_3$};
  \draw[flow] (a1)--(m1);
  \draw[flow] (a2)--(m2);
  \draw[flow] (a3)--(m3);
  \node[product, right=14mm of m2] (K) {$\mat{K}=\displaystyle\sum_t\mat{A}_t$};
  \draw[flow] (m1) -- (K);
  \draw[flow] (m2) -- (K);
  \draw[flow] (m3) -- (K);
  \node[font=\scriptsize\itshape, simGray, above=2mm of t1, xshift=6mm]
    {list of weak-form terms};
\end{tikzpicture}
\caption[The term fold $\mat{K}=\sum_t\mat{A}_t$]{The outer sum of
\cref{eq:termsum}. An analysis specifies a list of weak-form terms; each term
$t=(Q,\diffop)$ assembles (via \cref{eq:assembly}, an inner sum over elements)
into its own global operator $\mat{A}_t$; the global operator of the analysis is
the sum $\mat{K}=\sum_t\mat{A}_t$. Adding a physical effect---a loss term, a
second material, the $i\omega\mat{C}$ damping---means appending one term to the
list, never rebuilding from scratch.}
\label{fig:termfold}
\end{figure}

\subsection{Naming the pattern: \texttt{foldr}}

\Cref{eq:termsum} is not merely a formula; it is a \emph{computational pattern}
with a name. We have a list, a way to turn each list element into a matrix, and a
way to combine matrices---add them, starting from the zero matrix. Walking a list
and combining its elements into a single accumulated result is the job of the
higher-order function \texttt{fold}. In the pseudo-language of Part~III, the
assembly operator \code{fe\_assemble} is written exactly this way:
\begin{lstlisting}[style=l4style]
fe_assemble :: (space: FiniteElementSpace, terms: [WeakFormTerm]) -> LinOp
fe_assemble space terms = foldr (\t acc -> assembleTerm space t + acc) zero terms
                        -- = sum over t in terms of  assembleTerm space t
\end{lstlisting}
Read the body left to right. \code{foldr} walks the list \code{terms} from the
right; for each term \code{t} it computes that term's operator
\code{assembleTerm\ space\ t} and \emph{adds} it to the accumulator \code{acc};
the walk starts from \code{zero}, the zero operator. Unfolded, the call is
\[
\texttt{assembleTerm}\,t_1+\bigl(\texttt{assembleTerm}\,t_2+\cdots
+(\texttt{assembleTerm}\,t_m+\mathbf 0)\bigr)
=\sum_t\mat{A}_t,
\]
which is \cref{eq:termsum} exactly. The single per-term map
\code{assembleTerm\ space\ t}---the realization of one $(Q,\diffop)$ term as a
global matrix---is itself the inner element-fold \cref{eq:assembly}; we treat it
here as an opaque leaf and fold over it. The point is the \emph{outer} structure:
assembly of an analysis is one fold over the term list.

\begin{keyidea}
\textbf{Assembly is a commutative-monoid fold over the list of weak-form terms.}
The ingredients are a combine operation (operator addition $+$), an identity for
it (the zero operator $\mathbf 0$), and a list (the weak-form terms). Operator
addition is associative and commutative and has $\mathbf 0$ as identity---it forms
a \emph{commutative monoid}---and folding a list with a commutative-monoid
operation is the cleanest combinator pattern there is. Three consequences follow
purely from this algebra, with no reference to any code: assembling no terms
yields the zero operator; assembling a concatenated list equals assembling the
parts and adding; and the order of the terms never matters. We prove all three in
\cref{sec:laws}. This is the reader's first deep encounter with a thesis Part~III
makes central---\emph{the simulator's structure is folds over lists}
(\cref{ch:fold}, \cref{ch:combinators}).
\end{keyidea}

\section{The three laws of the assembly fold}
\label{sec:laws}

The commutative-monoid structure is not decoration; it is a set of guarantees an
engineer relies on. We state them as a theorem and prove them from the algebra of
operator addition alone---so they hold for \emph{every} analysis, regardless of
which terms are in the list or how the per-term operators are computed.

\begin{theorem}[Algebraic laws of assembly]
\label{thm:foldlaws}
Let $\mathrm{asm}(\mathit{terms})=\sum_{t\in\mathit{terms}}\mat{A}_t$ be the
assembled global operator on a fixed finite-element space, where $\mat{A}_t$ is
the per-term operator and $+$ is operator addition with identity the zero operator
$\mathbf 0$. Then:
\begin{enumerate}[label=\textup{(\roman*)},nosep]
\item \emph{(Empty-list identity)} $\mathrm{asm}([\,])=\mathbf 0$.
\item \emph{(Concatenation homomorphism)} for any two term lists $T_1,T_2$,
\[
\mathrm{asm}(T_1\!+\!\!+\,T_2)=\mathrm{asm}(T_1)+\mathrm{asm}(T_2),
\]
where $+\!\!+$ is list concatenation.
\item \emph{(Order independence)} $\mathrm{asm}(\mathit{terms})$ is unchanged
under any permutation of $\mathit{terms}$.
\end{enumerate}
\end{theorem}

\begin{proof}
All three are properties of the sum, inherited from operator addition.

\emph{(i)} The empty list has no summands, and the empty sum is by convention the
identity of the operation summed: $\sum_{t\in[\,]}\mat{A}_t=\mathbf 0$. (This is
the base case of the fold: \code{foldr\ f\ zero\ [\,]\ =\ zero}.)

\emph{(ii)} Concatenating two lists concatenates their summands, and a sum over a
concatenated index set splits:
\[
\mathrm{asm}(T_1\!+\!\!+\,T_2)
=\sum_{t\in T_1+\!\!+\,T_2}\mat{A}_t
=\sum_{t\in T_1}\mat{A}_t+\sum_{t\in T_2}\mat{A}_t
=\mathrm{asm}(T_1)+\mathrm{asm}(T_2).
\]
The middle equality is associativity of $+$ (regroup the summands into the two
sublists). A map from list-concatenation to a combine on results is exactly a
\emph{list homomorphism}; (ii) says assembly \emph{is} one.

\emph{(iii)} A permutation reorders the summands but changes neither the set of
summands nor their total, because $+$ is commutative and associative: any
reordering of $\mat{A}_{t_1}+\cdots+\mat{A}_{t_m}$ gives the same operator. The
list of terms is therefore effectively an unordered \emph{bag}.
\end{proof}

\begin{remark}
The three laws are not independent novelties; they are the universal property of
folding over a commutative monoid, specialized to (operators, $+$, $\mathbf 0$).
Any fold over any commutative monoid satisfies the analogues. We will meet the
same three laws again when we fold over \emph{elements} inside a single
$\mat{A}_t$ (the inner sum of \cref{eq:termsum}, also a commutative-monoid fold),
and yet again, abstracted away from finite elements entirely, in
\cref{ch:combinators}. Recognizing the pattern once buys it everywhere.
\end{remark}

\begin{workedexample}{The concatenation homomorphism, numerically}{concat-homo}
\label{wex:concat-homo}
Law (ii) of \cref{thm:foldlaws} says we may assemble a combined term list, or
assemble the pieces and add, and get the same operator. We check it on the
$3$-node 1-D mesh of \cref{wex:assemble-1d}, with a two-term list: a stiffness
term $t_K=(1,\Grad)$ and a mass term $t_M=(1,\mathrm{id})$.

\emph{The two per-term operators.} The stiffness operator is the
$\mat{K}=\tfrac1h\operatorname{tridiag}(-1,2,-1)$ of \cref{wex:assemble-1d}; with
$h=\tfrac14$, $\mat{K}=4\operatorname{tridiag}(-1,2,-1)$. The mass operator is the
$\mat{M}=\tfrac h6\operatorname{tridiag}(1,4,1)$ of \cref{wex:mass}; with
$h=\tfrac14$, $\mat{M}=\tfrac1{24}\operatorname{tridiag}(1,4,1)$. (Each is itself
an inner element-fold, assembled exactly as in \cref{wex:assemble-1d}.)

\emph{Route A: assemble the concatenated list.} Feed the assembler the single
list $[\,t_K,t_M\,]$. Its fold computes
$\mathrm{asm}([t_K,t_M])=\mat{A}_{t_K}+\mat{A}_{t_M}=\mat{K}+\mat{M}$:
\[
\mat{K}+\mat{M}
=4\!\begin{pmatrix}2&-1&0\\-1&2&-1\\0&-1&2\end{pmatrix}
+\frac1{24}\!\begin{pmatrix}4&1&0\\1&4&1\\0&1&4\end{pmatrix}
=\begin{pmatrix}
\tfrac{49}{6} & -\tfrac{95}{24} & 0\\[2pt]
-\tfrac{95}{24} & \tfrac{49}{6} & -\tfrac{95}{24}\\[2pt]
0 & -\tfrac{95}{24} & \tfrac{49}{6}
\end{pmatrix}.
\]
(Diagonal: $8+\tfrac{4}{24}=8+\tfrac16=\tfrac{49}{6}$. Off-diagonal:
$-4+\tfrac{1}{24}=-\tfrac{95}{24}$.)

\emph{Route B: assemble each, then add.} Call the assembler twice---once on
$[t_K]$, once on $[t_M]$---obtaining $\mat{K}$ and $\mat{M}$ separately, and add
the two matrices. The result is the very same $\mat{K}+\mat{M}$ above, entry for
entry. The two routes agree, as law (ii) promised: $\mathrm{asm}([t_K,t_M])
=\mathrm{asm}([t_K])+\mathrm{asm}([t_M])$.

\emph{Why it matters.} Route B is what an analysis actually does when it caches
the stiffness once and reuses it while sweeping a frequency-dependent combination
$\mat{K}-\omega^2\mat{M}$ (the driven analysis, \cref{ch:driven}): law (ii)
guarantees the cached-and-combined operator is bit-for-bit the from-scratch
assembly. It is also the law that makes assembly \emph{splittable across
processors}---hand half the term list (or half the elements) to each of two
machines, assemble independently, and add the results.
\end{workedexample}

\section{What is \emph{not} in the fold}
\label{sec:notinfold}

A fold's discipline is as much about what it excludes as what it includes. Two
operations that beginners often imagine are part of assembly are deliberately
\emph{outside} it.

\subsection{Boundary conditions are separable post-compositions}

The assembled operator $\mat{K}=\sum_t\mat{A}_t$ of \cref{eq:termsum} knows
nothing about boundary conditions. Essential (Dirichlet) conditions---pinning
certain degrees of freedom to prescribed values---are imposed \emph{after} the
fold, by acting on the already-assembled $\mat{K}$: the rows and columns of the
pinned dofs are eliminated, and any inhomogeneous data is lifted into the
right-hand side. Crucially, this elimination is a function of the assembled
operator alone; it does not care which terms were folded to build it, or in what
order. Assembly builds the operator; boundary-condition elimination is a
\emph{separable post-composition} on it, valid independently of the assembly. This
is why \cref{thm:foldlaws} can hold without any boundary-condition hypothesis: the
laws are about the fold, and the fold ends before the boundary conditions begin.
\Cref{ch:bc} takes up the elimination in full, including why keeping it separate
preserves symmetry and how inhomogeneous data moves to the load vector.

\begin{pitfall}
It is tempting to ``apply the boundary conditions as you assemble''---skipping the
pinned rows inside the element loop. This works, but it fuses two stages that are
cleaner apart, and it breaks the fold laws: an assembler that eliminates
mid-fold is no longer a homomorphism, because $\mathrm{asm}(T_1\!+\!\!+\,T_2)$
would eliminate twice. Keep the fold pure and post-compose the elimination once;
\cref{thm:foldlaws} is the reward. We assembled the \emph{interior} dofs in
\cref{wex:assemble-1d} precisely by dropping boundary slots, but notice we dropped
them from the \emph{connectivity} ($\mat{R}_K$), not by altering the fold's combine
step---the fold itself stayed a clean sum.
\end{pitfall}

\subsection{The matrix need not be formed at all}

The second exclusion is subtler and points forward to \cref{ch:matrixfree}.
Equation~\cref{eq:assembly} \emph{forms} the global matrix $\mat{K}$---it
materializes every nonzero entry into storage. But the only thing the iterative
solvers of \cref{part:framework} ever ask of $\mat{K}$ is its \emph{action} on a
vector, $\vw\mapsto\mat{K}\vw$. And the action distributes through both folds
without ever building the matrix:
\begin{equation}
\mat{K}\vw=\Bigl(\sum_t\sum_{K}\mat{R}_K^{\mathsf T}\mat{A}^{\mathrm{elem}}_{t,K}\mat{R}_K\Bigr)\vw
=\sum_t\sum_{K}\mat{R}_K^{\mathsf T}\bigl(\mat{A}^{\mathrm{elem}}_{t,K}(\mat{R}_K\vw)\bigr).
\label{eq:matrixfree-preview}
\end{equation}
Read the right-hand side as an algorithm: for each term and each element, gather
the element's slice of $\vw$ (apply $\mat{R}_K$), hit it with the small element
matrix, scatter-add the result (apply $\mat{R}_K^{\mathsf T}$); sum over all
elements and terms. No global matrix is ever assembled---only the small element
matrices (or, in libCEED, not even those) and the gather/scatter tables are kept.
This is the \emph{matrix-free} operator, and it is exactly the same fold, with the
combine step ``add into the global matrix'' replaced by ``add into the global
\emph{output vector}.''~\Cref{ch:matrixfree} develops it; \cref{eq:matrixfree-preview}
is the seed.

There is one more layer to peel, which we record now because it justifies the
notation $\mat{R}_K=G$ and points squarely at the production code. The element
matrix $\mat{A}^{\mathrm{elem}}_{t,K}$ is itself never formed in a high-performance
matrix-free kernel; it factors as
\begin{equation}
\mat{A}^{\mathrm{elem}}_{t,K}=B^{\mathsf T}_{\diffop}\,D\,B_{\diffop},
\label{eq:bdb}
\end{equation}
where $B_{\diffop}$ \emph{evaluates} the differential operator $\diffop$ of the
term at the element's quadrature points (the ``basis apply''), $D$ is a diagonal
of \emph{quadrature weights times the material coefficient $Q$} at those points,
and $B^{\mathsf T}_{\diffop}$ integrates back. Substituting \cref{eq:bdb} into
\cref{eq:matrixfree-preview}, one term's global action is the five-stage pipeline
\begin{equation}
\mat{A}_t=\;G^{\mathsf T}\;B^{\mathsf T}_{\diffop}\;D\;B_{\diffop}\;G,
\label{eq:gtbdbg}
\end{equation}
read right to left: gather to elements ($G$), evaluate at quadrature points
($B_{\diffop}$), scale by weighted material values ($D$), integrate back
($B^{\mathsf T}_{\diffop}$), scatter-add to global ($G^{\mathsf T}$). Only the
diagonal $D$ carries any floating-point ``physics''; $G$ and $B$ are fixed,
problem-independent index-and-interpolation maps. This factorization is the
contract of the libCEED library Palace assembles through, and its vocabulary---an
``$E$-vector'' is the element-local $G\vw$, an ``$L$-vector'' is the global $\vw$,
a ``$Q$-vector'' is the quadrature-point data---is precisely the gather/scatter
language of \cref{sec:gatherscatter}. We will not need \cref{eq:gtbdbg} again until
\cref{ch:matrixfree}, but it is worth seeing once that the entire matrix-free
apply is \emph{still the same two folds}, with the per-element work itself expanded
into a tiny pipeline.

\section{Sparsity, storage, and cost}
\label{sec:cost}

We close with the practical consequences of the fold structure: where the
sparsity comes from, what it costs to assemble, and where the parallelism lives.

\subsection{Why the assembled matrix is sparse}

The sparsity of $\mat{K}$ that \cref{ch:galerkin} attributed to the local support
of the basis is, from the assembly viewpoint, even more transparent. Each term
$\mat{R}_K^{\mathsf T}\mat{A}^{\mathrm{elem}}_{t,K}\mat{R}_K$ in \cref{eq:assembly}
is nonzero \emph{only} on the $L\times L$ block of rows and columns indexed by
element $K$'s degrees of freedom. A global entry $\mat{K}_{ij}$ can be nonzero
only if some single element $K$ owns \emph{both} dof $i$ and dof $j$---i.e.\ only
if nodes $i$ and $j$ are vertices of a common cell. Since each node belongs to a
bounded number of cells (two in 1-D, a small fan in 2-D/3-D), each row of
$\mat{K}$ has only a bounded number of nonzeros, independent of how large the mesh
grows. The matrix has $\mathcal{O}(N)$ nonzeros, not $N^2$. The sparsity
\emph{pattern} is read straight off the gather tables: $\mat{K}_{ij}\neq0$ iff
$i,j$ share an element. \Cref{fig:sparsitypattern} shows a small 2-D example.

\begin{figure}[t]
\centering
\begin{tikzpicture}[scale=1.0, every node/.style={font=\scriptsize}]
  % left: tiny 2D mesh (3x3 nodes, split into triangles)
  \begin{scope}[shift={(0,0)}]
    \foreach \x in {0,1,2}\foreach \y in {0,1,2}{
      \fill[simInk] (\x,\y) circle (1.4pt);}
    % grid lines + diagonals
    \foreach \y in {0,1,2}{\draw[simGray] (0,\y)--(2,\y);}
    \foreach \x in {0,1,2}{\draw[simGray] (\x,0)--(\x,2);}
    \foreach \x in {0,1}\foreach \y in {0,1}{\draw[simGray] (\x,\y+1)--(\x+1,\y);}
    % label center node 4 and highlight its star
    \node[simRust, below right, font=\tiny] at (1,1) {$4$};
    \fill[simRust] (1,1) circle (2.0pt);
    \node[font=\tiny] at (1,-0.5) {mesh: node $4$ couples to its neighbours};
  \end{scope}
  % right: the sparsity grid (9x9) with node 4's row marked
  \begin{scope}[shift={(4.6,-0.3)}, scale=0.34]
    \foreach \r in {0,...,8}\foreach \c in {0,...,8}{
      \draw[simGray!35] (\c,-\r) rectangle ++(0.94,-0.94);}
    % node numbering: bottom-left=0..8 row-major; node4=center couples to 1,3,4,5,7 (and diag 0,8 via triangles)
    \def\nbrs{1,3,4,5,7,0,8}
    \foreach \c in {1,3,4,5,7,0,8}{
      \fill[simBlue!45] (\c,-4) rectangle ++(0.94,-0.94);
      \fill[simBlue!45] (4,-\c) rectangle ++(0.94,-0.94);}
    \fill[simRust!60] (4,-4) rectangle ++(0.94,-0.94);
    \draw[simInk,thick] (0,0) rectangle (9,-9);
    \node[anchor=west, scale=2.2, align=left] at (9.6,-4.5) {row/col\\of node $4$};
  \end{scope}
\end{tikzpicture}
\caption[Sparsity from local support]{Sparsity of the assembled matrix on a small
$3\times3$-node triangular mesh. \emph{Left:} the centre node $4$ shares an
element with only its immediate mesh neighbours. \emph{Right:} the resulting
sparsity pattern. Row $4$ of $\mat{K}$ (and, by symmetry, column $4$) has a
nonzero only in the columns of those neighbours; all other entries of the row are
exactly zero, because no single element owns both node $4$ and a far node. Each
row has a bounded number of nonzeros no matter how large the mesh grows: $\mat{K}$
has $\mathcal{O}(N)$ nonzeros, and is stored as a sparse matrix---never as the
full $N\times N$ array.}
\label{fig:sparsitypattern}
\end{figure}

\subsection{The cost of assembly, and the scatter-add race}

The element-fold runs over $n_e$ elements; on each it computes a dense
$L\times L$ element matrix (a fixed amount of work set by the element type and
polynomial degree, not by $N$) and scatters it. Assembly therefore costs
\begin{equation}
\mathcal{O}\bigl(n_e\times(\text{per-element work})\bigr)
=\mathcal{O}(N)
\label{eq:asmcost}
\end{equation}
total work, linear in the problem size---one of the cheapest stages of the whole
pipeline. And because the element-fold is a commutative-monoid fold
(\cref{thm:foldlaws}(iii)), the elements may be processed in \emph{any} order,
and---more powerfully---in \emph{parallel}: order independence is exactly the
licence to split the list across processors. Assembly is, in this sense,
embarrassingly parallel over elements.

\begin{figure}[t]
\centering
\begin{tikzpicture}[scale=1.0, every node/.style={font=\footnotesize}]
  % four element "workers" computing in parallel
  \foreach \i/\x/\col in {1/0/simBlue,2/2.2/simRust,3/4.4/simGreen,4/6.6/simViolet}{
    \node[opbox, draw=\col, minimum width=14mm] (w\i) at (\x,1.5)
       {worker\\$K_\i$};
  }
  \node[font=\scriptsize\itshape, simGray] at (3.3,2.4) {parallel: element-local work is independent};
  % global vector slots
  \foreach \k/\lab in {0/1,1/2,2/3,3/4,4/5}{
    \draw[simInk] (1.6+\k*0.9,-0.6) rectangle ++(0.8,0.8);
    \node[font=\scriptsize] at (2.0+\k*0.9,-0.2) {$x_{\lab}$};
  }
  \node[anchor=east, font=\footnotesize\bfseries] at (1.5,-0.2) {global $\vx$:};
  % scatter-add arrows: K1->x1,x2 ; K2->x2,x3 (collision on x2) etc
  \draw[flow, simBlue] (w1) -- (2.0,0.25);
  \draw[flow, simBlue] (w1) -- (2.9,0.25);
  \draw[flow, simRust] (w2) -- (2.9,0.25);
  \draw[flow, simRust] (w2) -- (3.8,0.25);
  \draw[flow, simGreen] (w3) -- (3.8,0.25);
  \draw[flow, simGreen] (w3) -- (4.7,0.25);
  \draw[flow, simViolet] (w4) -- (4.7,0.25);
  \draw[flow, simViolet] (w4) -- (5.6,0.25);
  % collision markers on shared dofs x2,x3,x4
  \foreach \k in {2,3,4}{
    \node[circle,draw=simRust,fill=simBgRust,inner sep=1pt,font=\tiny]
      at (1.1+\k*0.9,0.55) {$\scriptstyle\rightleftarrows$};}
  \node[simRust!80!simInk, font=\scriptsize, align=left, anchor=west] at (6.2,0.5)
     {shared dof:\\two writers\\$\Rightarrow$ race};
\end{tikzpicture}
\caption[Parallel over elements, race on shared dofs]{Assembly parallelizes over
elements---each worker computes its element's contribution independently
(top)---but the \emph{scatter-add} into the global vector is where independence
breaks down. A degree of freedom shared by two elements has two workers trying to
add into the same slot at once (the $\rightleftarrows$ markers). Because the
underlying operation is addition, the \emph{result} is order-independent
(\cref{thm:foldlaws}(iii)); but two simultaneous read--modify--write updates can
interleave and lose one, a \emph{data race}. The race is resolved with atomic
adds, per-colour scheduling, or thread-private buffers summed at the end---an
implementation concern that does not change the assembled value.}
\label{fig:parallelrace}
\end{figure}

There is exactly one place where the parallelism is not free, and the fold laws
pinpoint it. The element-local \emph{computation} is perfectly independent across
elements; the \emph{scatter-add} is not, because two elements sharing a degree of
freedom write into the same global slot (\cref{fig:parallelrace}). The
\emph{mathematical} result is order-independent---that is precisely
\cref{thm:foldlaws}(iii)---but two threads performing
read--modify--write on the same memory location can interleave and lose an
update: a \emph{data race}. The fix is a matter of implementation, not
mathematics: make the additions atomic, colour the elements so that no two
sharing a dof run simultaneously, or give each thread a private accumulator and
sum them at the end (itself one more application of \cref{thm:foldlaws}(ii)).
Whichever is chosen, the assembled operator is the same: the race threatens the
\emph{schedule}, never the \emph{answer}, because the answer is a commutative-
monoid fold.

\begin{physicsbox}
The two parallel structures of assembly mirror the physics. Element-local work is
parallel because, in the continuum, the contribution of each little region of
space to the weak form is independent---fields do not interact across a region
boundary except through the shared boundary itself. The scatter-add ``race'' on
shared dofs is the discrete echo of that one exception: a node on the interface
between two elements feels both, and its equation must collect both
contributions. Superposition of local energies is the physics; scatter-add is its
arithmetic; the commutative-monoid fold is its structure.
\end{physicsbox}

\summarysection
Global assembly turns small, dense element matrices into the one large sparse
global operator the solver sees, by a pair of indexing operations. The
\emph{gather} $\mat{R}_K$ copies an element's global degrees of freedom into local
order; the \emph{scatter-add} $\mat{R}_K^{\mathsf T}$ writes element-local results
back, \emph{accumulating} where elements share a degree of freedom. The global
operator is the element sum $\mat{K}=\sum_K\mat{R}_K^{\mathsf
T}\mat{K}^{\mathrm{elem}}_K\mat{R}_K$, and the diagonal $\sum_K\mat{R}_K^{\mathsf
T}\mat{R}_K$ counts each dof's element-multiplicity. An analysis assembles a
\emph{list of weak-form terms}, each term assembling (an inner element-sum) into a
per-term operator $\mat{A}_t$, and the global operator is the outer sum
$\mat{K}=\sum_t\mat{A}_t$. Both sums are folds: assembly is a
\emph{commutative-monoid fold}, \code{fe\_assemble} $=$ \code{foldr (+) zero},
whose three laws---empty-list identity, concatenation homomorphism, and order
independence---follow from operator addition alone and underwrite caching,
splitting, and parallelism. Boundary-condition elimination is a \emph{separable
post-composition} on the assembled operator (\cref{ch:bc}), not part of the fold;
and the matrix need not be formed at all---the same fold, scattering into an
output vector instead of a matrix, is the matrix-free operator of
\cref{ch:matrixfree}. Sparsity follows from local support (a row is nonzero only
in the columns sharing an element), assembly costs $\mathcal{O}(N)$, and it is
embarrassingly parallel over elements save for the scatter-add race on shared
dofs---a race in the \emph{schedule}, never the \emph{answer}. That assembly is a
fold over a list is the reader's first deep meeting with the thread Part~III makes
the spine of the whole book.

\furtherreading
The element-by-element assembly procedure, local-to-global maps, and the sparse
storage of the resulting matrix are covered in every finite-element text; Brenner
and Scott~\citep{brennerscott2008} give the mathematical framing (the assembly
operator as a sum of restrictions) while Jin~\citep{jin2014} gives the engineering
practice for electromagnetics specifically, including the element-matrix
computations for the $\Hcurl$ and $\Hdiv$ bases we have only previewed. The
gather/scatter-and-element-matrix decomposition $A=G^{\mathsf
T}B^{\mathsf T}DBG$, and the matrix-free assembly that exploits it, are the
organizing idea of the libCEED project~\citep{ceed2021}, whose ``$E$-vector /
$L$-vector'' vocabulary is precisely the element-local versus global distinction
of \cref{sec:gatherscatter}; readers headed for \cref{ch:matrixfree} will find the
CEED interface a faithful realization of \cref{eq:matrixfree-preview}. The
recognition of assembly as a fold and a list homomorphism, and the parallel
decompositions it licenses, are developed abstractly in \cref{ch:combinators}.

\exercisessection
\begin{enumerate}[nosep]
  \item \textbf{Multiplicity by hand.} For the four-element 1-D mesh of
  \cref{wex:assemble-1d}, write out the diagonal matrix $\sum_K
  \mat{R}_K^{\mathsf T}\mat{R}_K$ of \cref{eq:multiplicity} on the three interior
  dofs. Confirm each entry is $2$ and explain, in one sentence, why the (omitted)
  boundary nodes would have multiplicity $1$.

  \item \textbf{Five elements.} Repeat \cref{wex:assemble-1d} for a mesh of $n=5$
  cells ($h=\tfrac15$, four interior dofs). List the five elements and the global
  indices of their local slots, then scatter-add the five $2\times2$ blocks and
  read off the $4\times4$ tridiagonal $\mat{K}$. Which dofs have multiplicity $2$?

  \item \textbf{A triangle's element matrix.} A single linear triangle has three
  local dofs (its vertices) and a $3\times3$ element matrix. Without computing its
  entries, write the assembly contribution $\mat{R}_K^{\mathsf
  T}\mat{K}^{\mathrm{elem}}\mat{R}_K$ for a triangle whose vertices are global dofs
  $\{2,5,7\}$: which $9$ entries of the global matrix does it touch, and what is
  the multiplicity contribution to dof $5$ if three triangles meet there?

  \item \textbf{Empty and singleton.} Using \cref{thm:foldlaws}, evaluate
  $\mathrm{asm}([\,])$ and $\mathrm{asm}([t])$ for a single term $t$. Then show
  $\mathrm{asm}([t_1,t_2])=\mathrm{asm}([t_2,t_1])$ directly from the two laws you
  need, naming each.

  \item \textbf{The homomorphism, again.} An eigenmode analysis needs both a
  stiffness term $t_K$ and a mass term $t_M$. Using law (ii), express
  $\mathrm{asm}([t_K,t_M])$ in terms of the separately-assembled $\mat{K}$ and
  $\mat{M}$. Then explain why the driven analysis, which forms
  $\mat{K}-\omega^2\mat{M}$ at many frequencies $\omega$, assembles $\mat{K}$ and
  $\mat{M}$ \emph{once} rather than re-running the fold at each frequency---and
  which law guarantees this is correct.

  \item \textbf{Matrix-free apply.} Starting from \cref{eq:matrixfree-preview},
  write out $\mat{K}\vw$ for the single-term four-element mesh of
  \cref{wex:assemble-1d} acting on $\vw=(1,0,0)^{\mathsf T}$, by gathering,
  applying the $2\times2$ element matrix, and scatter-adding---without ever forming
  $\mat{K}$. Check your answer against the first column of the tridiagonal
  $\mat{K}$.

  \item \textbf{Boundary conditions outside the fold.} Suppose you ``optimize'' the
  assembler to skip the equation of a pinned dof \emph{inside} the element loop,
  before summing. Give a two-term list $T_1\!+\!\!+\,T_2$ for which this fused
  assembler violates the concatenation homomorphism (ii), and identify the double-
  counting. (Hint: what happens to the pinned dof's row when both sublists pass
  through the elimination?)

  \item \textbf{The scatter-add race.} On the parallel picture of
  \cref{fig:parallelrace}, suppose two threads each execute
  \texttt{x[2] = x[2] + (local contribution)} on the shared dof $x_2$ without
  synchronization. Describe a specific interleaving of their reads and writes that
  loses one contribution, and explain why \cref{thm:foldlaws}(iii) nonetheless
  guarantees the \emph{intended} result is independent of which thread goes first.
  Which of the three fixes named in the text (atomics, colouring, private buffers)
  relies on law (ii)?
\end{enumerate}
